\reading{MR-15}{The Art of Term Structure Models: Volatility and Distribution}{strike}{7}
% ==========================================================

\begin{mrkeybox}[title={Learning objectives in this reading}]
\small
\begin{enumerate}[label=\textbf{\alph*.},leftmargin=1.7em]
\item Describe the short-term rate process under a model with time-dependent volatility.
\item Calculate the short-term rate change and determine the behavior of the standard deviation
      of the rate change using a model with time-dependent volatility.
\item Assess the efficacy of time-dependent volatility models.
\item Describe the short-term rate process under the Cox-Ingersoll-Ross (CIR) and lognormal
      models.
\item Calculate the short-term rate change and describe the basis point volatility using the CIR
      and lognormal models.
\item Describe lognormal models with deterministic drift and mean reversion.
\end{enumerate}
\end{mrkeybox}

% ---------------------------------------------------------
\lo{MR-15 a}{Describe the short-term rate process under a model with time-dependent volatility.}

This model extends previous rate models by incorporating \term{time-dependent volatility}. It is
a natural extension of the Ho-Lee model, which includes non-constant volatility.

\begin{mrfmlbox}
\[ dr \;=\; \underbrace{\lambda(t)\,dt}_{\text{time-dependent \textbf{drift}}}
   \;+\; \underbrace{\sigma(t)\,dw}_{\text{time-dependent \textbf{volatility}}} \]
\mrvk{The drift term is the one carried over from the previous reading. The volatility term is
what is new here: $\sigma$ is now a function of $t$ rather than a constant.}{}
\end{mrfmlbox}

% ---------------------------------------------------------
\lo{MR-15 b}{Calculate the short-term rate change and determine the behavior of the standard deviation of the rate change using a model with time-dependent volatility.}

\subsection{Model 3 --- time-dependent drift and volatility}
Model 3 is a special case of time-dependent volatility models. It augments the Ho-Lee model with
time-dependent volatility: the volatility of the short rate starts at the constant $\sigma$ and
\term{exponentially declines to zero}.

\begin{mrfmlbox}
\[ dr \;=\; \lambda(t)\,dt \;+\; \sigma e^{-\alpha t}\,dw \]
\mrvk{$\sigma$ = volatility at $t = 0$, which decreases exponentially to 0 for $\alpha > 0$;\quad
$\alpha$ = the decay rate;\quad $\lambda(t)$ = the time-dependent drift.}{}
\end{mrfmlbox}

\begin{mrexbox}[title={Example --- short-term rate change under Model 3}]
Initial short-term rate $r_0 = 5\%$, drift $\lambda = 0.24\%$, time-dependent volatility
$\sigma e^{-0.3t}$ starting at 1.50\%. The random variable $dw$ from a normal distribution is
0.2, with mean 0 and standard deviation $\sqrt{1/12} = 0.2887$.
\tcblower
\small
\begin{align*}
dr &= 0.24\% \times \left(\tfrac{1}{12}\right) \;+\; 1.5\% \times e^{-0.3(1/12)} \times 0.2\\
dr &= 0.02\% + 0.29\% \;=\; \mathbf{0.31\%}
\end{align*}
\end{mrexbox}

\subsection{The standard deviation of the terminal distribution}

\begin{center}
\begin{tikzpicture}
\begin{axis}[
  width=10.8cm, height=5.4cm,
  xmin=0, xmax=31, ymin=0, ymax=6.2,
  xlabel={\scriptsize Horizon (years)},
  ylabel={\scriptsize Standard deviation (\%)},
  xlabel style={font=\scriptsize}, ylabel style={font=\scriptsize},
  tick label style={font=\scriptsize},
  xtick={0,5,10,15,20,25,30}, ytick={0,1,2,3,4,5,6},
  axis line style={draw=black!55}, clip=false]
\addplot[FRMGold, dashed, line width=0.8pt, forget plot]
  coordinates {(0,5.635)(31,5.635)};
\node[anchor=east, font=\scriptsize\sffamily, text=FRMGold, fill=white, inner sep=1.5pt]
  at (axis cs:30.5,5.95) {asymptote $= \sigma/\sqrt{2\alpha} = 5.63\%$};
\addplot[FRMNavy, line width=1.3pt, smooth] coordinates {
(0,0)(1,1.244)(2,1.738)(3,2.103)(5,2.650)(7,3.062)(10,3.535)(13,3.896)
(15,4.093)(18,4.341)(20,4.480)(23,4.658)(25,4.760)(28,4.891)(30,4.967)};
\end{axis}
\end{tikzpicture}
\end{center}
\figcap{Standard deviation of the terminal distribution of short rates in Model 3, with
$\sigma = 126$ bp and $\alpha = 0.025$. Because volatility decays, the cumulative uncertainty
stops growing and flattens toward a ceiling rather than rising without bound as $\sqrt{T}$ would
in a constant-volatility model.}

% ---------------------------------------------------------
\lo{MR-15 c}{Assess the efficacy of time-dependent volatility models.}

\begin{mrkeybox}[title={Model 3 effectiveness}]
\begin{enumerate}
\item \term{More flexibility} for the short-term rate.
\item Useful for pricing \term{multi-period derivatives} like interest rate caps and floors.
\item \term{Volatility term structure:}
  \begin{enumerate}[label=\alph*)]
  \item If volatility starts at a constant $\sigma$ and then exponentially declines to zero,
        volatility cannot be forecast for the longer term, so the volatility term structure is
        \textbf{declining}.
  \item Volatility is determined \textbf{independently} of the level of interest rates.
  \end{enumerate}
\item It leads to \term{parallel shifts} in the yield curve.
\item \term{Model usage:}
  \begin{enumerate}[label=\alph*)]
  \item Valuations and hedging: use Vasicek or a mean reversion model.
  \item Pricing options on fixed income products: use a Model 3 type, which accounts for
        volatility changes.
  \end{enumerate}
\end{enumerate}
\end{mrkeybox}

\begin{mrtrapbox}[title={Model 3 versus Vasicek --- the sibling pair}]
\begin{itemize}
\item The behaviour of the standard deviation as a function of the time horizon in Model 3
      \textit{resembles} the effect of mean reversion.
\item If the initial volatility and decay rate $\alpha$ in Model 3 are set equal to the
      volatility and mean reversion rate $k$ in the Vasicek model, the \term{terminal standard
      deviations for both models become identical}.
\item \term{Key difference:} Model 3 is a \textbf{parallel} shift model with no mean reversion.
      Vasicek is a \textbf{nonparallel} shift model, because of the mean reversion effect.
\end{itemize}
That is the trap: identical terminal standard deviations, completely different shift behaviour.
\end{mrtrapbox}

% ---------------------------------------------------------
\lo{MR-15 d}{Describe the short-term rate process under the Cox-Ingersoll-Ross (CIR) and lognormal models.}
\lo{MR-15 e}{Calculate the short-term rate change and describe the basis point volatility using the CIR and lognormal models.}

\begin{mrtrapbox}[title={Issues with the existing models}]
\term{Basis point volatility.} Existing models determine the basis point volatility of the
short-term rate \textit{without considering the level of the short-term rate itself}. Two key
concerns follow:
\begin{enumerate}
\item \term{Extreme inflation.} During periods of hyperinflation, expected rate changes for the
      next period are substantially larger than in normal conditions.
\item \term{Ultra-low interest rates.} In an environment with near-zero or slightly negative
      interest rates, rate volatility naturally diminishes.
\end{enumerate}
\end{mrtrapbox}

\subsection{Solution 1 --- the Cox-Ingersoll-Ross model}
CIR accounts for the fact that the volatility of interest rates can change as a function of the
level of the short-term rate, and that interest rates appear to have a mean-reverting feature.

\begin{mrfmlbox}
\[ dr \;=\; k(\theta - r)\,dt \;+\; \sigma\sqrt{r}\;dw \]
\mrvk{$k$, $\theta$ = the mean reversion speed and long-run level, as in Vasicek;\quad
$\sigma\sqrt{r}$ = the basis point volatility, which now \textit{rises with the level of the
rate}. Volatility increases by the \textbf{square root} of the rate, so the increase is
non-linear.}{}
\end{mrfmlbox}

\begin{mrexbox}[title={Example --- CIR short-term rate calculation}]
Current short-term rate 5\%, long-run value 24\%, mean reversion speed $k = 0.04$, volatility
$\sigma = 1.50\%$, and a $dw$ realization of 0.2. Find the change in the short-term rate after
one month.
\tcblower
\small
\begin{align*}
dr &= 0.04\,(24\% - 5\%)\left(\tfrac{1}{12}\right) \;+\; 1.5\% \sqrt{5\%} \times 0.2\\
dr &= 0.063\% + 0.067\% \;=\; \mathbf{0.13\%}
\end{align*}
\end{mrexbox}

\begin{mrtrapbox}[title={The unit trap in $\sqrt{r}$}]
The rate inside the square root is taken as a \textbf{decimal}, not as a percentage number.
Here $\sqrt{5\%} = \sqrt{0.05} = 0.2236$, giving $1.5\% \times 0.2236 \times 0.2 = 0.067\%$.
Using $\sqrt{5} = 2.236$ instead would give 0.67\%, ten times too large, and that wrong answer
is exactly the kind of distractor that gets offered.
\end{mrtrapbox}

\subsection{Solution 2 --- lognormal models (Model 4)}
The lognormal model maintains \term{constant yield volatility}, but basis point volatility will
\term{increase with the short-term rate level}.

\begin{mrfmlbox}
\[ dr \;=\; a\,r\,dt \;+\; \sigma\,r\,dw \]
\mrvk{$\sigma r$ = the basis point volatility, proportional to the level of the rate. Volatility
increases \textbf{linearly} with the rate, in contrast to CIR's square root.}{}
\end{mrfmlbox}

\subsection{How the three volatility specifications compare}

\begin{center}
\begin{tikzpicture}
\begin{axis}[
  width=10.8cm, height=5.6cm,
  xmin=0, xmax=20, ymin=0, ymax=2.3,
  xlabel={\scriptsize Level of the short-term rate, $r$ (\%)},
  ylabel={\scriptsize Basis point volatility (indexed to 1.0 at $r = 10\%$)},
  xlabel style={font=\scriptsize}, ylabel style={font=\scriptsize},
  tick label style={font=\scriptsize},
  xtick={0,5,10,15,20}, ytick={0,0.5,1.0,1.5,2.0},
  legend style={font=\scriptsize\sffamily, draw=FRMRule, fill=white,
    at={(0.5,-0.28)}, anchor=north, legend columns=-1, column sep=10pt},
  axis line style={draw=black!55}, clip=false]
\addplot[black!65, line width=1.2pt, dashed, domain=0:20, samples=2] {1};
\addlegendentry{Normal models --- $\sigma$ constant}
\addplot[FRMNavy, line width=1.3pt, smooth, domain=0:20, samples=120] {sqrt(x/10)};
\addlegendentry{CIR --- $\sigma\sqrt{r}$}
\addplot[FRMRed, line width=1.3pt, dotted, domain=0:20, samples=2] {x/10};
\addlegendentry{Lognormal --- $\sigma r$}
\draw[black!35, line width=0.5pt, dashed] (axis cs:10,0) -- (axis cs:10,1);
\node[circle, fill=black!55, inner sep=1.5pt] at (axis cs:10,1) {};
\end{axis}
\end{tikzpicture}
\end{center}
\figcap{All three curves are set equal at a 10\% rate so only the \textit{shape} differs. A
normal model gives the same basis point volatility whether rates are at 1\% or 19\%. CIR bends
over, because the square root grows more slowly than the rate. The lognormal model rises in a
straight line through the origin, so at zero rates it gives zero volatility.}

\subsection{Comparison between CIR and lognormal (terminal distributions)}

\begin{center}
\small
\begin{tabularx}{\linewidth}{@{}p{0.6cm} >{\raggedright\arraybackslash}X@{}}
\arrayrulecolor{FRMRule}\toprule
\rowcolor{FRMNavy} \hdr{} & \hdr{Point}\\
a) & Both CIR and lognormal are designed to ensure the short-term rate remains
  \term{non-negative}, a significant advantage when modeling interest rates. Both have a
  \term{positive drift}, preventing negative rates.\\
\rowcolor{FRMSoft}
b) & In both cases \term{yield-based volatility is constant}; only basis point volatility
  changes.\\
c) & The \term{normal} distribution implies a positive probability of \textit{negative} interest
  rates, especially over longer forecast horizons, which can be a drawback. CIR is more complex
  but might be a better fit if volatility's dependence on the rate is crucial.\\
\rowcolor{FRMSoft}
d) & Both result in terminal distributions that are always \term{non-negative and
  right-skewed}. This has important implications for pricing out-of-the-money options and other
  instruments, where the distribution shape can impact pricing.\\
\bottomrule
\end{tabularx}
\end{center}

\begin{center}
\begin{tikzpicture}
\begin{axis}[
  width=10.8cm, height=5.0cm,
  xmin=-4, xmax=16, ymin=0, ymax=0.30,
  xlabel={\scriptsize Terminal short rate (\%)},
  axis y line=none, axis x line=bottom, axis line style={draw=black!55},
  xlabel style={font=\scriptsize}, tick label style={font=\scriptsize},
  xtick={-4,0,4,8,12,16},
  legend style={font=\scriptsize\sffamily, draw=FRMRule, fill=white,
    at={(0.5,-0.30)}, anchor=north, legend columns=-1, column sep=10pt},
  clip=false]
\addplot[domain=-4:0, samples=60, smooth, draw=none, fill=FRMRed!30, forget plot]
  {0.1329*exp(-((x-6)^2)/18)} \closedcycle;
\addplot[black!70, line width=1.2pt, dashed, smooth, domain=-4:16, samples=160]
  {0.1329*exp(-((x-6)^2)/18)};
\addlegendentry{Normal --- symmetric, mass below zero}
\addplot[FRMNavy, line width=1.3pt, smooth, domain=0.05:16, samples=160]
  {0.62*exp(-((ln(x)-1.55)^2)/0.42)/x};
\addlegendentry{CIR / lognormal --- non-negative, right-skewed}
\draw[FRMRed, line width=0.9pt] (axis cs:0,0) -- (axis cs:0,0.25);
\node[anchor=west, font=\scriptsize\sffamily, text=FRMRed, fill=white, inner sep=1.5pt,
  align=left] at (axis cs:-3.8,0.215) {the normal model puts\\ real probability here};
\draw[FRMRed, -{Stealth[length=3.5pt]}, line width=0.6pt]
  (axis cs:-1.9,0.205) -- (axis cs:-1.4,0.135);
\end{axis}
\end{tikzpicture}
\end{center}
\figcap{The shaded sliver to the left of zero is the whole problem with a normal model: it
assigns genuine probability to negative rates, and that mass grows with the forecast horizon.
CIR and lognormal terminal distributions stop dead at zero and lean right, which is what makes
them better behaved for pricing out-of-the-money options.}

% ---------------------------------------------------------
\lo{MR-15 f}{Describe lognormal models with deterministic drift and mean reversion.}

\subsection{1. The Salomon Brothers model --- lognormal with deterministic drift}
In this model the natural logarithm of the short-term rate follows a normal distribution, and it
is used to construct an interest rate tree based on that natural logarithm.

\begin{mrfmlbox}
\[ d\bigl[\ln(r)\bigr] \;=\; a(t)\,dt \;+\; \sigma\,dw \]
\mrvk{$a(t)$ = the drift term, which varies across time just as in the Ho-Lee model. The crucial
change is that the process is written on $\ln(r)$, not on $r$.}{}
\end{mrfmlbox}

\begin{center}
\begin{tikzpicture}[every node/.style={font=\tiny\sffamily},
  lg/.style={rectangle, rounded corners=2pt, draw=FRMNavy!45, fill=PaleNavy,
    line width=0.6pt, inner sep=3.5pt, align=center},
  rt/.style={rectangle, rounded corners=2pt, draw=FRMGreen!50, fill=PaleGreen,
    line width=0.6pt, inner sep=3.5pt, align=center}]
% log space
\begin{scope}
  \node[lg] (a) at (0,0) {$\ln r_0$};
  \node[lg] (u) at (2.5,1.1) {$\ln r_0 + a_1 dt + \sigma\sqrt{dt}$};
  \node[lg] (d) at (2.5,-1.1) {$\ln r_0 + a_1 dt - \sigma\sqrt{dt}$};
  \node[lg] (uu) at (6.4,2.0) {$\ln r_0 + (a_1{+}a_2)dt + 2\sigma\sqrt{dt}$};
  \node[lg] (mm) at (6.4,0) {$\ln r_0 + (a_1{+}a_2)dt$};
  \node[lg] (dd) at (6.4,-2.0) {$\ln r_0 + (a_1{+}a_2)dt - 2\sigma\sqrt{dt}$};
  \foreach \p/\q in {a/u, a/d, u/uu, u/mm, d/mm, d/dd}{
    \draw[FRMNavy!55, line width=0.6pt] (\p) -- (\q);}
  \node[anchor=north, text=black!60] at (3.2,-2.8) {\textbf{the tree is built in $\ln(r)$ space}};
\end{scope}
\end{tikzpicture}
\end{center}

\begin{center}
\begin{tikzpicture}[every node/.style={font=\tiny\sffamily},
  rt/.style={rectangle, rounded corners=2pt, draw=FRMGreen!50, fill=PaleGreen,
    line width=0.6pt, inner sep=3.5pt, align=center}]
  \node[rt] (a) at (0,0) {$r_0$};
  \node[rt] (u) at (2.6,1.1) {$r_0\,e^{(a_1 dt + \sigma\sqrt{dt})}$};
  \node[rt] (d) at (2.6,-1.1) {$r_0\,e^{(a_1 dt - \sigma\sqrt{dt})}$};
  \node[rt] (uu) at (6.6,2.0) {$r_0\,e^{((a_1+a_2)dt + 2\sigma\sqrt{dt})}$};
  \node[rt] (mm) at (6.6,0) {$r_0\,e^{(a_1+a_2)dt}$};
  \node[rt] (dd) at (6.6,-2.0) {$r_0\,e^{((a_1+a_2)dt - 2\sigma\sqrt{dt})}$};
  \foreach \p/\q in {a/u, a/d, u/uu, u/mm, d/mm, d/dd}{
    \draw[FRMGreen!55, line width=0.6pt] (\p) -- (\q);}
  \node[anchor=north, text=black!60] at (3.3,-2.8)
    {\textbf{exponentiating gives the rate tree}};
\end{tikzpicture}
\end{center}
\figcap{The same tree twice. Above, additive steps in $\ln(r)$. Below, the same tree
exponentiated back into rates, where the steps become \textit{multiplicative}.}

\begin{mrkeybox}[title={Why this guarantees positive rates}]
\begin{itemize}
\item Future movements of the short rate in a lognormal model are \term{multiplicative} instead
      of being \term{additive} as in normal models and the Ho-Lee model.
\item The implication is that all rates in the tree will \textbf{always be positive}, since
      $e^{x} > 0$ for all $x$.
\end{itemize}
\end{mrkeybox}

\subsection{2. Lognormal model with mean reversion --- Black-Karasinski}
The lognormal distribution combined with a mean-reverting process is known as the
\term{Black-Karasinski} model. This model is very flexible, allowing for time-varying volatility
\textit{and} mean reversion.

\begin{mrfmlbox}
\[ d\bigl[\ln(r)\bigr] \;=\; k(t)\bigl[\ln\theta(t) - \ln(r)\bigr]\,dt \;+\; \sigma(t)\,dw \]
\mrvk{$k(t)$ = time-varying mean reversion speed;\quad $\theta(t)$ = time-varying long-run
level;\quad $\sigma(t)$ = time-varying volatility. Every one of the three parameters can move
through time, which is what makes the model flexible.}{}
\end{mrfmlbox}

\begin{mrtrapbox}
This model does \term{not} produce a naturally recombining interest rate tree. To force the tree
to recombine, the time steps $dt$ must be \term{recalibrated}.
\end{mrtrapbox}

% ---------------------------------------------------------
\clearpage
\section{Summary of all the term structure models}

The source notes close MR-14 and MR-15 with a combined summary. It is reproduced here as a
single reference table.

\begin{center}
\small
\begin{tabularx}{\linewidth}{@{}p{2.6cm} p{5.0cm} >{\raggedright\arraybackslash}X@{}}
\arrayrulecolor{FRMRule}\toprule
\rowcolor{FRMNavy} \hdr{Model} & \hdr{Process} & \hdr{What it adds}\\
\textbf{Model 1}
  & $dr = \sigma\,dw$
  & No drift; rates normally distributed.\\
\rowcolor{FRMSoft}
\textbf{Model 2}
  & $dr = \lambda\,dt + \sigma\,dw$
  & Adds a positive drift to Model 1, interpretable as a risk premium for longer horizons.\\
\textbf{Ho-Lee}
  & $dr = \lambda(t)\,dt + \sigma\,dw$
  & Generalises the drift to be time-dependent.\\
\rowcolor{FRMSoft}
\textbf{Vasicek}
  & $dr = k(\theta - r)\,dt + \sigma\,dw$
  & Makes the drift depend on the current rate: mean reversion.\\
\textbf{Model 3}
  & $dr = \lambda(t)\,dt + \sigma e^{-\alpha t}\,dw$
  & Time-dependent \textit{volatility}: $\sigma$ at $t=0$, decaying exponentially to 0 for
    $\alpha > 0$.\\
\rowcolor{FRMSoft}
\textbf{CIR}
  & $dr = k(\theta - r)\,dt + \sigma\sqrt{r}\,dw$
  & Mean reverting with constant $\sigma$ but basis point volatility $\sigma\sqrt{r}$, which
    increases at a \textit{decreasing} rate.\\
\textbf{Model 4}\newline (lognormal, deterministic drift)
  & $d[\ln(r)] = a(t)\,dt + \sigma\,dw$
  & Yield volatility $\sigma$ constant, but basis point volatility $\sigma r$ increases with the
    level of the rate.\\
\rowcolor{FRMSoft}
\textbf{Model 4}\newline (lognormal, mean reverting)
  & $d[\ln(r)] = k(t)[\ln\theta(t) - \ln(r)]\,dt + \sigma(t)\,dw$
  & Black-Karasinski. Adds mean reversion and time-varying volatility in log space.\\
\bottomrule
\end{tabularx}
\end{center}

\begin{mrkeybox}[title={The one-line way to hold all eight}]
Read down the process column and watch what changes. Models 1, 2 and Ho-Lee vary the
\textbf{drift constant}. Vasicek makes the drift depend on \textbf{where the rate is}. Model 3
moves to varying the \textbf{volatility over time}. CIR and the lognormal models make the
volatility depend on \textbf{the level of the rate} --- by $\sqrt{r}$ and by $r$ respectively.
Black-Karasinski does both at once.
\end{mrkeybox}


% ==========================================================
